\section*{Hoofdstuk \ref{ch:b3:atmospheric-dynamics} --- Atmosferische dynamica}

\begin{solution}{exo:b3:atmospheric-dynamics:1}
(a)~$\zeta = 2\omega = 2\times2\pi/86\,400 = \qty{1.45e-4}{s^{-1}}$.
(b)~$f = \num{7.29e-5}$ en \qty{1.26e-4}{s^{-1}}: de \omterm{def:b3:atmospheric-dynamics:eddy}{wervel} heeft ongeveer
tweemaal de \omterm{def:b3:atmospheric-dynamics:vorticity-kinds}{planetaire vorticiteit} van \ang{30}~NB en iets meer dan die van
\ang{60}~NB.
\end{solution}

\begin{solution}{exo:b3:atmospheric-dynamics:2}
$\zeta = f_{30} - f_{40} = \num{7.29e-5} - \num{9.37e-5} =
\qty{-2.08e-5}{s^{-1}}$: anticyclonisch.
\end{solution}

\begin{solution}{exo:b3:atmospheric-dynamics:3}
$k = 2\pi/(5\times10^6) = \qty{1.257e-6}{m^{-1}}$; $c = -\beta/k^2 =
-\num{1.62e-11}/\num{1.58e-12} = \qty{-10.3}{m/s}$: westwaarts.
\end{solution}

\begin{solution}{exo:b3:atmospheric-dynamics:4}
$\beta = 2\Omega\cos\ang{60}/a = \qty{1.14e-11}{m^{-1}.s^{-1}}$;
$\lambda_s = 2\pi\sqrt{10/\num{1.14e-11}} = \qty{5870}{km}$; de cirkel,
\qty{20015}{km}, bevat $n_s = 3.4$ golven.
\end{solution}

\begin{solution}{exo:b3:atmospheric-dynamics:5}
$q = f/h$ blijft behouden: over het gebergte $\zeta = f(8/10) - f =
\qty{-2.06e-5}{s^{-1}}$, anticyclonisch; over het bekken $f(12/10) - f =
\qty{+2.06e-5}{s^{-1}}$, cyclonisch. Waar kolommen worden uitgerekt, in de luwte
van het gebergte en boven het bekken, ontstaat cyclonische vorticiteit:
daar zijn depressies geneigd zich te vormen.
\end{solution}

\begin{solution}{exo:b3:atmospheric-dynamics:6}
$k = 5/(a\cos\ang{45}) = \qty{1.11e-6}{m^{-1}}$, $\beta/k^2 =
\qty{13.1}{m/s}$: $c = \qty{1.9}{m/s}$, $c_g = \qty{28.1}{m/s}$. De energie
legt \qty{5000}{km} af in \num{2.1}~dagen, het patroon in \num{31}~dagen.
\end{solution}

\begin{solution}{exo:b3:atmospheric-dynamics:7}
(a)~$k^2 = \beta/U - l^2 = \num{1.079e-12} - \num{2.74e-13} =
\qty{8.05e-13}{m^{-2}}$, $k = \qty{8.97e-7}{m^{-1}}$: $\lambda =
\qty{7000}{km}$ en $n = 4.0$, minder golven dan bij $l = 0$. (b)~$l^2 =
\num{1.097e-12} > \beta/U = \num{1.079e-12}$: $k^2$ zou negatief zijn.
\end{solution}

\begin{solution}{exo:b3:atmospheric-dynamics:8}
$L_d = 0.012\times10^4/\num{1.031e-4} = \qty{1164}{km}$; $\sigma =
0.31\times\num{1.031e-4}\times\num{2e-3}/0.012 = \qty{5.3e-6}{s^{-1}}$, een
$e$-vouwtijd van \num{2.2}~dagen; golflengte $2\pi\times1164/1.61 =
\qty{4540}{km}$.
\end{solution}

\begin{solution}{exo:b3:atmospheric-dynamics:9}
$[v^*T^*] = (VT_0/2)\cos\delta$: \qty{25}{K.m/s}, \qty{12.5}{K.m/s} en 0;
warmtefluxen \num{25100}, \num{12600} en \qty{0}{W/m^2}. Wanneer de
temperatuur een kwart golflengte uit fase is met de wind, gaat evenveel warme
lucht naar het zuiden als naar het noorden: geen transport.
\end{solution}

\begin{solution}{exo:b3:atmospheric-dynamics:10}
Linearisatie rond $U$ met $\zeta' = \nabla^2\psi'$ en $v' =
\partial_x\psi'$: $(\partial_t + U\partial_x)\nabla^2\psi' + \beta\partial_x
\psi' = 0$. Met $\psi' \propto \mathrm e^{\mathrm i(kx + ly - \omega t)}$,
volgt $(\omega - Uk)K^2 = -\beta k$, dus $\omega = Uk - \beta k/K^2$. Ten
opzichte van de lucht is de fasesnelheid $c - U = -\beta/K^2 < 0$ voor elke golf.
\end{solution}

\begin{solution}{exo:b3:atmospheric-dynamics:11}
$\varphi_H = [5\times9.81\times12\,000\times0.25/(3\times(\num{7.292e-5})^2
\times(\num{6.371e6})^2)]^{1/2} = \qty{0.477}{rad}$, \ang{27.3}.
$u_M = \Omega a\varphi_H^2 = \qty{106}{m/s}$; exact geldt $\Omega a\sin^2
\varphi_H/\cos\varphi_H = \qty{110}{m/s}$: tweeënhalf keer de
waargenomen straalstroom.
\end{solution}

\begin{solution}{exo:b3:atmospheric-dynamics:12}
(a)~Over de band is het grensvlak verplaatst met $\eta = sy$, $-L/2 \le y
\le L/2$, ten opzichte van de vlakke ligging. Een laag met verminderd dichtheidscontrast
die over $\eta$ is verplaatst, slaat $\tfrac12\rho g'\eta^2$ per eenheid oppervlak op; gemiddeld
over de band geldt $\langle\eta^2\rangle = s^2L^2/12$. (b)~$g' = 0.327$:
\num{1.5e5} en \qty{6.1e5}{J/m^2}, tegenover $\tfrac12\times0.7\times15^2
\times10^4 = \qty{7.9e5}{J/m^2}$ aan kinetische energie: vergelijkbaar.
\end{solution}

\begin{solution}{pb:b3:atmospheric-dynamics:1}
\textbf{1.} $f_{45} = \qty{1.031e-4}{s^{-1}}$, $f_{50} =
\qty{1.117e-4}{s^{-1}}$; $\zeta = f_{45} - f_{50} = \qty{-8.6e-6}{s^{-1}}$.

\textbf{2.} $R = V/|\zeta| = 15/\num{8.6e-6} = \qty{1750}{km}$; met de klok mee,
terug naar het zuiden.

\textbf{3.} $\zeta = f_{45}(9.5/10) - f_{50} = \qty{-1.38e-5}{s^{-1}}$, meer
anticyclonisch.

\textbf{4.} $\zeta = f_{45} - f_{40} = \qty{+9.4e-6}{s^{-1}}$, cyclonisch: het
draait weer naar het noorden. Hoe het pakketje ook verplaatst is, het krijgt de
vorticiteit die het terugdraait: een herstellend effect, en een oscillatie.

\textbf{5.} $\beta = \qty{1.62e-11}{m^{-1}.s^{-1}}$; de cirkel is
\qty{28306}{km} lang.

\textbf{6.} $n = 3$ tot 8: golflengten 9435, 7076, 5661, 4718, 4044 en
\qty{3538}{km}; $c = -21.5$, $-5.5$, 1.9, 5.9, 8.3 en \qty{9.9}{m/s}; $c_g
= 51.5$, 35.5, 28.1, 24.1, 21.7 en \qty{20.1}{m/s}.

\textbf{7.} $n = 3$ en 4 bewegen westwaarts, $n \ge 5$ oostwaarts. De golf $n = 5$
doet $10^6/1.86 = \qty{5.4e5}{s}$, \num{6.2}~dagen, over \qty{1000}{km}.

\textbf{8.} Stroomafwaarts, oostwaarts: bij $c_g = \qty{24.1}{m/s}$ legt de energie
\qty{5000}{km} af in \num{2.4}~dagen.

\textbf{9.} $\lambda_s = 2\pi\sqrt{15/\num{1.62e-11}} = \qty{6050}{km}$,
$n_s = 4.7$.

\textbf{10.} $n_s = 5.7$ voor \qty{10}{m/s} en 4.1 voor \qty{20}{m/s}: een
sterkere westenwind begunstigt langere stationaire golven, minder rond
het halfrond.

\textbf{11.} $n_s = 5.8$ bij \ang{35} NB en 3.4 bij \ang{55} NB.

\textbf{12.} $\zeta = f(8/10) - f = \qty{-2.06e-5}{s^{-1}}$: samendrukking boven
de bergen maakt de kolom anticyclonisch, uitrekking in de lij
maakt haar cyclonisch. Het gebergte legt de stroming een vast patroon van vorticiteit op,
een forcering die niet beweegt.

\textbf{13.} Een vaste forcering geeft een stationaire respons; de respons is
het sterkst wanneer het opgelegde patroon zelf een vrije golf is die stil zou
staan, de resonantie bij $\lambda_s$.

\textbf{14.} $\Lambda = 25/10^4 = \qty{2.5e-3}{s^{-1}}$;
$|\partial T/\partial y| = \Lambda fT/g = \num{2.5e-3}\times\num{1.031e-4}
\times250/9.81 = \qty{6.6}{K}$ per \qty{1000}{km}.

\textbf{15.} $\mathrm{Ri} = 0.011^2/(\num{2.5e-3})^2 = 19$: ruim boven 0.25,
geen schuifturbulentie.

\textbf{16.} $L_d = 0.011\times10^4/\num{1.031e-4} = \qty{1067}{km}$.

\textbf{17.} $2\pi\times1067/1.61 = \qty{4160}{km}$; $28306/4160 = 6.8$
golven.

\textbf{18.} $2\pi\times1067/2.40 = \qty{2790}{km}$.

\textbf{19.} $\sigma = 0.31\times\num{1.031e-4}\times\num{2.5e-3}/0.011 =
\qty{7.3e-6}{s^{-1}}$, \num{0.63}~dag$^{-1}$.

\textbf{20.} $\exp(3\times0.627) = 6.6$.

\textbf{21.} De storm, \qty{4160}{km}, is korter dan de \omterm{prop:b3:atmospheric-dynamics:stationary}{stationaire
Rossby-golflengte}, \qty{6050}{km}: als Rossby-golf wordt hij door de westenwind
oostwaarts meegevoerd, en de stormen trekken oostwaarts door de langzamere of
stilstaande lange golven.

\textbf{22.} Met de hoogte naar het westen gekanteld: alleen dan is de naar de pool
stromende lucht warmer dan de naar de evenaar stromende lucht, zodat warmte
poolwaarts gaat en de \omterm{def:b3:atmospheric-dynamics:ape}{beschikbare potentiële energie} vrijkomt.

\textbf{23.} De snelst groeiende storing groeit met een factor $e$ elke
\textbf{\num[text-series-to-math=true]{1.6}~dagen} (\num{1.59}~dagen volgens de
berekening).
\end{solution}
